\chapterimage{images/02-poslednje-hodocasce.jpeg}
\chapter{Poslednje hodočašće}

\editorialnote{Sedam hodočasnika, njihovi prvi utisci, pravila hodočašća i odluka da ispričaju svoje priče.}

Konzul se probudio iz kriogene fuge nekoliko sati pre dolaska na Hiperion.

Glava ga je bolela, grlo mu je bilo suvo, a misli usporene. Pored njega su stajala dva člana posade i visoki čovek u tamnoj odori sa kapuljačom.

Bio je to Het Mastin, kapetan Igdrasila i templarski Pravi Glas Drveta.

``Još pet sati do Hiperiona'', rekao mu je. ``Ostali su već budni. Čekaju vas.''

Igdrasil nije ličio ni na jedan brod kojim je Konzul ranije putovao.

Bio je doslovno ogromno živo drvo, gotovo kilometar dugo, prilagođeno putovanju između zvezda. Deblo, grane i lišće bili su isprepletani sa motorima, poljima sile, kabinama i platformama. Na granama su se nalazili prolazi, mostovi i vrtovi, a između listova su svetlucale veštačke i prave zvezde.

Templari su takve brodove smatrali živim bićima.

Konzul ih je smatrao jednim od retkih čuda Hegemonije koja još nisu izgubila sposobnost da ga iznenade.

Dok su išli prema platformi za ručavanje, primetio je nešto neobično.

Brod je bio gotovo prazan.

``Gde su putnici?'' pitao je.

``Vi ste jedini.''

Konzul se zaustavio.

Drvobrod poput Igdrasila obično je prevozio hiljade ljudi. Putovanje bez putnika značilo je ogroman gubitak.

Onda je shvatio.

Igdrasil nije došao samo zbog hodočasnika.

Ako počne rat, mogao je da posluži za evakuaciju.

To je značilo da Hegemonija zaista očekuje najgore.

Na otvorenoj platformi, visoko među granama, čekalo ga je petoro ljudi.

Het Mastin je seo na čelo stola.

Konzul je zauzeo prazno mesto pored njega i počeo da posmatra ostale.

Prvi je bio otac Lenar Hojt.

Katolički sveštenik, mlad po godinama, ali ne i po izgledu. Bio je mršav, iscrpljen i očigledno bolestan. Lice mu je izgledalo kao lice čoveka koji već dugo živi sa bolom.

Konzulu je njegovo ime nešto značilo.

Pre mnogo godina na Hiperionu je nestao drugi katolički sveštenik.

Veza mu se nije dopadala.

Do Hojta je sedeo pukovnik Fedman Kasad.

Njega nije trebalo predstavljati.

Kasad je bio poznati vojnik Hegemonije, heroj za jedne i ratni zločinac za druge. Bio je visok, vitak i potpuno miran. Nije mnogo govorio, ali Konzul je odmah imao osećaj da taj čovek primećuje sve oko sebe.

Preko puta njega sedeo je Martin Silenus.

Pesnik.

Nizak, bučan, vulgaran i očigledno već pripit.

Delovao je kao čovek koji se trudi da svima oko sebe bude neprijatan i pritom iskreno uživa u tome.

Konzul nije bio siguran da li mu je zabavan ili nepodnošljiv.

Verovatno oba.

Pored Silenusa sedeo je Sol Vejntraub.

Stariji naučnik kratke sede brade i umornih očiju.

Za razliku od ostalih, u naručju je držao bebu.

Konzul je znao ko je dete.

Rejčel.

Čuo je ponešto o Vejntraubovoj priči i njegovoj neobičnoj potrazi, ali nije očekivao da će dete biti toliko malo.

Poslednja među njima bila je Bron Lamija.

Privatna detektivka sa Lususa.

Bila je niska i snažna, naviknuta na gravitaciju mnogo veću od zemaljske. Posmatrala je Konzula otvoreno i procenjivački, bez pokušaja da sakrije šta radi.

Nije mu se učinila kao osoba kojoj bi bilo lako lagati.

Konzul je prebrojao prisutne.

Hojt.

Kasad.

Silenus.

Vejntraub i dete.

Lamija.

On.

Šest odraslih.

Okrenuo se ka Mastinu.

``Rekli ste da nas je sedmoro.''

Mastin je klimnuo.

``Jeste.''

Konzul je pogledao Rejčel.

``Dete?''

``Ne. Hodočasnik mora svesno da odluči da potraži Šrajka.''

Nastala je kratka tišina.

``Ja sam sedmi'', rekao je Het Mastin.

To je promenilo raspoloženje za stolom.

Templar nije bio samo njihov prevoznik.

I on je išao sa njima do kraja.

Sedam hodočasnika.

Crkva Šrajka je vekovima slala grupe prema Vremenskim grobnicama, ali ovo hodočašće bilo je drugačije.

Možda poslednje.

U sistemu se već osećalo prisustvo rata.

Proterani još nisu stigli, rekao im je Mastin, ali njihovi brodovi bili su svega nekoliko dana iza njih.

``Hoće li biti rata?'' pitao je Hojt.

Niko nije odmah odgovorio.

Konzul je konačno slegnuo ramenima.

``Ne znam. Niko zapravo ne zna šta Proterani žele.''

Kasad se nije složio sa nadom da bi njihov dolazak mogao da prođe mirno.

Ako osvoje Hiperion, rekao je, neće ga jednostavno okupirati. Rat sa Proteranima nije bio rat između država koje žele nove granice. Ako odluče da unište planetu, mogu da je ostave potpuno mrtvu.

Posle toga niko nekoliko trenutaka nije govorio.

Hegemonija ih je pratila ratnim brodom.

Kada je Konzul kroz vojni dvogled video koliko ozbiljno naoružan brod leti uz Igdrasil, postalo mu je još jasnije koliko je situacija opasna.

Nije ih čuvala obična patrola.

Flota se pripremala za bitku.

Ipak, hodočašće se nastavljalo.

``Zašto jednostavno ne sletimo kod Vremenskih grobnica?'' pitao je Hojt. ``Ako već dolazi rat, zašto gubimo nekoliko dana na putovanje preko planete?''

Het Mastin je odmahnuo glavom.

Hodočašće je imalo svoja pravila.

Počinjalo je u prestonici, a završavalo se kod Šrajka.

``Pravila'', promrmljala je Lamija. ``Usred rata.''

Konzul je razumeo njenu frustraciju, ali je znao i razlog zbog kog niko nije pokušavao da zaobiđe tradiciju.

Ljudi su već pokušavali da priđu Vremenskim grobnicama letelicama.

Brodovi bi sleteli.

Kasnije bi se vratili.

Na njima ne bi bilo nikoga.

Nije bilo tragova borbe.

Nije bilo oštećenja.

Nije bilo neobičnih energetskih pražnjenja, prinudnog menjanja kursa ili znakova da je bilo šta pošlo naopako.

Samo bi ljudi nestali.

``Dakle, brodovi se vraćaju'', rekla je Lamija.

``Da.''

``A putnici?''

Konzul je odmahnuo glavom.

``Ne.''

Het Mastin je mirno dodao:

``Nikada.''

Silenus se nasmejao.

``Odlično. Rat iza nas, čudovište ispred nas i niko nema pojma zašto smo ovde. Lepo putovanje.''

To je, pomislio je Konzul, bio neprijatno precizan opis njihove situacije.

Tokom večere razgovor je postao nešto ličniji.

Sol Vejntraub prvi je izgovorio pitanje koje je očigledno mučilo sve za stolom.

``Da li iko od vas zna zašto je izabran?''

Niko nije odgovorio.

``Dobro'', rekao je Sol. ``Onda drugo pitanje. Da li je iko od nas uopšte pripadnik Crkve Šrajka?''

Ponovo gotovo niko.

Sol je bio Jevrejin.

Konzul ateista.

Hojt katolički sveštenik.

Lamija nije imala nikakav interes za religiju.

Kasad nije ponudio jasan odgovor, ali ničim nije pokazivao pripadnost kultu.

Silenus je tvrdio da je tokom dugog života promenio toliko vera i filozofija da ih više ni sam verovatno nije mogao nabrojati.

Na kraju je sebe nazvao paganom.

``A za jednog pagana'', rekao je uz osmeh, ``Šrajk je sasvim pristojan bog.''

Het Mastin jeste bio templar, ali ni on nije pripadao Crkvi Šrajka. Neki templari videli su Šrajka kao neku vrstu kazne ili božanskog bića, ali Mastin je takvo verovanje smatrao jeresi.

Sol je pogledao oko stola.

``Zar vam to nije neobično?''

Milioni ljudi želeli su da krenu na hodočašće.

Crkva je mogla da izabere svoje najvernije sledbenike.

Umesto toga izabrala je njih.

Ljude koji ne veruju u nju.

Ljude koji, svaki na svoj način, već imaju neku vezu sa Hiperionom.

``Možda upravo zato treba da razgovaramo'', rekao je Sol.

Kasad ga je pogledao.

``O čemu?''

``O nama.''

Lamija se namrštila.

``Zašto?''

Sol je spustio pogled ka Rejčel, koja mu je spavala na grudima.

``Zato što svako od nas zna nešto što ostali ne znaju.''

Objasnio je da međuzvezdana putovanja nisu samo razdvojila ljude prostorom.

Razdvojila su ih i vremenom.

Zbog vremenskog duga neko je mogao da ima šezdeset godina, a da je rođen pre više od jednog veka. Ljudi za istim stolom mogli su da pripadaju različitim periodima istorije, iako su fizički bili približno vršnjaci.

Svako od njih video je drugi deo Hiperiona.

Drugi deo njegove prošlosti.

Možda je svaki nosio deo iste slagalice.

``Ako već idemo prema nečemu što nas verovatno želi mrtve'', rekao je Sol, ``voleo bih makar da razumemo zašto.''

``I šta predlažete?'' pitao je Konzul.

``Da svako ispriča svoju priču.''

Kasad nije video smisao.

Lamija je pitala šta ih sprečava da lažu.

``Ništa'', odgovorio je Silenus. ``Zato će biti zanimljivo.''

Konzul je ćutao.

Njega je zanimalo nešto drugo.

Meina Gledston mu je rekla da se među sedmoricom nalazi agent Proteranih.

Ako budu govorili o sebi, možda će neko napraviti grešku.

Možda će priče otkriti više nego što njihovi pripovedači žele.

``Glasajmo'', rekao je.

Kasad ga je pogledao sa blagim podsmehom.

``Od kada smo demokratija?''

``Od kada moramo zajedno da stignemo do Šrajka.''

Predlog je prihvaćen većinom glasova.

Svi će ispričati zbog čega su povezani sa Hiperionom.

Ostalo je samo pitanje redosleda.

Silenus je iscepao komad papira na sedam delova, napisao brojeve i ubacio ih u Konzulov šešir.

Izvlačili su jedan po jedan.

Konzul je otvorio svoju cedulju tako da je niko ne vidi.

Broj sedam.

Osetio je olakšanje.

Do tada je moglo da se dogodi bilo šta.

Mogli su da stignu do Šrajka.

Mogao je da počne rat.

Mastin je mogao da promeni plan.

Grupa je mogla da odustane od priča.

A možda niko od njih neće ni doživeti sedmi dan.

``Ko je prvi?'' pitao je Silenus.

Otac Lenar Hojt podigao je svoju cedulju.

Na njoj je stajao broj jedan.

Nije izgledao srećno zbog toga.

``Dobro'', rekao je Silenus. ``Počnite.''

Hojt je nekoliko trenutaka ćutao.

Zatim je ustao i otišao.

Kada se vratio, u rukama je nosio dve stare, nagorele sveske.

Spustio ih je na sto.

``Ako treba da vam ispričam svoju priču'', rekao je, ``moraću prvo da vam ispričam priču jednog drugog čoveka.''

Konzul je pogledao dnevnike.

Nije znao šta se u njima nalazi.

Ali znao je da je upravo počelo ono zbog čega će sedam stranaca, pre nego što stignu do Šrajka, morati da otkriju jedni drugima stvari koje možda nikada nikome nisu nameravali da kažu.

Prvi je bio sveštenik.
